% !TEX program = xelatex
\documentclass[lang=zh]{akizumi-book}
\akizumisetup{kicker={國立臺灣大學\akizumisep 碩士論文}}
\title{自我預測表徵的分佈式評論家}
\author{Aaron Lin}

\begin{document}
\frontmatter
\maketitle
\tableofcontents

\mainmatter
\chapter{緒論}\label{ch:intro}
\chaptertoc
\section{動機}
模型為本的表徵方法在 2026 年大量出現，例如 MR.Q 與 LeWM。
\begin{definition}[Collapse]\label{def:collapse}
若 $\Var[f(\vb{o})] = 0$，稱編碼器 $f$ collapse。
\end{definition}
\section{本文貢獻}
見\cref{ch:method}與\cref{def:collapse}。

\chapter{方法}\label{ch:method}
\chaptertoc
\section{問題設定}
\begin{equation}\label{eq:loss}
  \mathcal{L} = \E\norm{\hat{\vb{z}}_{t+1} - \sg(\vb{z}_{t+1})}^2 .
\end{equation}
\begin{theorem}\label{thm:a}
由\cref{eq:loss}可得結論。
\end{theorem}
\begin{proof} 顯然。 \end{proof}
\section{實驗}
\subsection{環境}
DMC 與 Atari。

\appendix
\chapter{證明細節}
補充\cref{thm:a}的證明。
\end{document}
